\documentclass[11pt]{article}

\input{0-parameters.tex}

\begin{document}

\section{Comments Emilien 14th April}
\subsection{Correcting typo}

There's a slight mistake in Eq. \ref{eq_dual_processed}, see below

\begin{align}
\hat{A}_i &= \er{-(}\hat{p}_i - \sum_j a_{ji} \hat{p}_{j}\er{)} \\
\hat{A}_i &= \er{-(}\frac{1}{p_i X_i} \left( (p_i X_i) \hat{p}_i - (p_i X_i) \sum_j a_{ji} \hat{p}_{j} \right)\er{)}\\
\hat{A}_i &= \er{-(}\frac{1}{p_i X_i} \left( (p_i X_i)\hat{p}_i - \sum_j X_{ji}\er{p_{j}} \hat{p}_{j}\right)\er{)}\\
\hat{A}_i &= \er{-(}\frac{1}{p_i X_i} \left( (p_i X_i)\frac{\Delta p_i}{p_i} - \sum_j X_{ji}p_{j} \frac{\Delta p_j}{p_j} \right)\er{)}\\
\hat{A}_i &= \er{-(}\frac{1}{p_i X_i} \left( X_i\Delta p_i - \sum_j X_{ji} \Delta p_j \right)\er{)}\\
\hat{A}_i &= \er{-(}\frac{1}{p_i X_i} \left( \Delta C_i  - \Delta C_i^{U} \right)\er{)} = \er{-(}\frac{\Delta C^*_i}{p_i X_i}\er{)},
\end{align}

Correcting it, we get
\begin{equation}
    \hat{A}_i  = -\frac{\Delta C^*_i}{p_i X_i}.
\end{equation}



\subsection{Calibrating $\zeta$ on the CCC scenario}

You show that by writing the dual we get for each sector that ``free of double counting'' costs correspond to TFP shocks, while ``raw'' costs correspond to price changes. Equating Eq. \eqref{eq_dual_processed} and Eq. \eqref{eq_ahat_t_tau_ccc}, we get
\begin{equation*}
    -\Delta C^*_i=(\zeta-1)\Delta C_i,
\end{equation*} 
hence 
\begin{equation}
\label{eq_zeta_ccc}
     \zeta=\frac{-\Delta C^*_i}{\Delta C_i}+1
\end{equation}

Let us look at that in the CCC data. In Table \ref{tab_TFP_annualised_nonprocessed_processed}, for each sector is the annualised resource cost (annualised processed cost) given by the CCC. My best effort at recomputing it (given that I don't have the discount rate nor lifespan). Last column is the annualised unprocessed cost using the same discount rate and lifespan as the annualised processed cost. The zeta columns correponds to Eq. \eqref{eq_zeta_ccc}.


\begin{longtable}{cP{3cm}P{3cm}P{3cm}P{3cm}}
\caption{Extra 2050 cost (bn£) per sector before/after
the re-allocation to remove double counting (Balanced Net
Zero, CCC) and compared to the CCC annualised resource cost - Annualised cost approach. \footnotesize \textit{Manual correction of the Fuel supply as the model does work great to annualise it for the non-processed cost. To be investigated}} \\ 
  \hline
Sector & CCC annual cost & Recomputed annual cost (Processed) & Recomputed annual cost (Non-processed) & zeta \\ 
  \hline
Agriculture & -0.99 & 0.18 & 0.19 & 0.08 \\ 
  Aviation & -1.27 & -0.93 & -0.86 & -0.08 \\ 
  Electricity supply & -5.30 & -2.62 & 1.40 & 2.88 \\ 
  F-gases & -0.00 & 0.01 & 0.01 & 0.00 \\ 
  Fuel supply & -1.56 & -1.69 & - & - \\ 
  LULUCF sinks & 2.62 & 2.53 & 2.53 & 0.00 \\ 
  LULUCF sources & -0.09 & 0.05 & 0.05 & 0.00 \\ 
  M\&C & 4.18 & 4.11 & 5.02 & 0.18 \\ 
  Non res buildings & 3.81 & 3.86 & 3.68 & -0.05 \\ 
  Removals & 5.71 & 6.36 & 6.57 & 0.03 \\ 
  Res buildings & 12.90 & 11.57 & 15.64 & 0.26 \\ 
  Shipping & 3.62 & 3.58 & 3.44 & -0.04 \\ 
  Surface transport & -9.09 & -7.79 & 1.29 & 7.03 \\ 
  Waste & 1.69 & 1.75 & 3.49 & 0.50 \\ 
   \hline
Total & 16.21 & 20.96 & 42.45 & 0.81 \\ 
  \hline
\label{tab_TFP_annualised_nonprocessed_processed}
\end{longtable}

Using the aggregate cost, we can compute productivity growth
\[
\hat{A}=-\frac{\sum_i \Delta C^*_i}{\text{GDP}}=-\frac{16.21}{2234}=-0.007
\]

With my own approximation of annualisation
\[
\hat{A}=-\frac{\sum_i \Delta C^*_i}{\text{GDP}}=-\frac{20.1}{2234}=-0.009
\]

With the non-processed costs, using  Eq. \eqref{eq_ahat_t_tau_ccc} and the calibrated $\zeta=0.13$
\[
\hat{A}=(\zeta-1)\frac{\sum_i \Delta C_i}{\text{GDP}}=(0.13-1)\frac{42.45}{2234}=-0.017
\]

We should have a $\zeta=\frac{-\sum_i \Delta C^*_i}{\sum_i \Delta C_i}+1=-\frac{20.1}{42.45}+1=0.52$ to account for the gap between processed and un-processed costs. 


\subsection{Do we just assume that $\zeta=0$ in the CCC application?}

We show a paradox when forcing $\hat{X}_M=0$ implying that $\zeta=0$.

As a reminder we consider this the simplified economy
\begin{itemize}
    \item \textbf{Upstream sector $U$} using only labour. 
    \[
    \hat{X}_U > 0,~
    \hat{p}_U < 0
    \]
    \item \textbf{Downstream sector $M$} using $U$ as its sole input. 
    \[
    \hat{X}_M = \hat{A}_M+a_{UM}  \hat{X}_U,~
    \hat{p}_M = -\hat{A}_M+a_{UM}  \hat{p}_U
    \]
\end{itemize}

% As soon as we assume that $\hat{X}_M=0$ and $\hat{X}_U=0$, then we get a problem. From the primal we get $\hat{X}_M = \hat{A}_M+a_{UM}  \hat{X}_U$, therefore, $\hat{A}_M=0$.

% It might not be a problem in a more complicated economy because we can enforce $\forall i, ~\hat{X}_i=0$ but still have $\hat{X}_{ji}\neq0$. Thus a reallocation of the production between sectors. 

We have two sectors, aggregate productivity is
\begin{equation*}
    \hat{A}=\lambda_M \hat{A}_M+\lambda_U \hat{A}_U.
\end{equation*}

As shown before, we consider only the downstream sector M, and show that we retrieve the true aggregate productivity growth.

\begin{equation}
    \begin{aligned}
    \hat{\tilde{A}}=\lambda_M  \hat{\tilde{A}}_M 
    &= \lambda_M\Bigl(\zeta \hat{X}_M + (\zeta-1) \hat{p}_M\Bigr)\\
    &= \lambda_M\Bigl(\zeta(\hat{A}_M+ \tfrac{\lambda_U}{\lambda_M} \hat{X}_U )
        +  (\zeta-1) (-\hat{A}_M+\tfrac{\lambda_U}{\lambda_M} \hat{p}_U)\Bigr)\\
    &= \lambda_M (\zeta\hat{A}_M-(\zeta-1)\hat{A}_M)+\lambda_U\Bigl(\zeta \hat{X}_U + (\zeta-1) \hat{p}_U\Bigr)\\
    &= \lambda_M \hat{A}_M+\lambda_U \hat{\tilde{A}}_U\\
    &=\lambda_M \hat{A}_M+\lambda_U \hat{A}_U\quad(\text{since we calibrate }\zeta\text{ such that }\hat{\tilde{A}}_U=\hat{A}_U)\\
    &=\hat{A}.
\end{aligned}
\end{equation}

Hence, equating the first and the last lines, we get
\begin{equation}
 \hat{A}=    \lambda_M\Bigl(\zeta \hat{X}_M + (\zeta-1) \hat{p}_M\Bigr).
\end{equation}

Which means, that in the particular case $\hat{X}_M=0$, we get $ \hat{A}=    \lambda_M(\zeta-1) \hat{p}_M.
$

But if I directly take into account that $\hat{X}_M=0$, we get
\begin{equation}
    \begin{aligned}
    \hat{\tilde{A}}=\lambda_M  \hat{\tilde{A}}_M 
    &= \lambda_M\Bigl((\zeta-1) \hat{p}_M\Bigr)\\
    &= \lambda_M(\zeta-1) (-\hat{A}_M+\tfrac{\lambda_U}{\lambda_M} \hat{p}_U)\\
    &= \lambda_M(1-\zeta) \hat{A}_M +\lambda_U(\zeta-1)\hat{p}_U
\end{aligned}
\end{equation}

Note that from the dual equation that $\hat{A}_U=-\hat{p}_U$, hence
\begin{equation} 
     \hat{\tilde{A}}= (1-\zeta)\left(\lambda_M \hat{A}_M +\lambda_U\hat{A}_U\right)=
     (1-\zeta)\hat{A}.
\end{equation}

These two approaches should give the same result, and they do only if $\zeta=0$. Hence, somewhere, we must do this assumption. 
In fact, for the upstream sector, we said that from the dual  $\hat{A}_U=-\hat{p}_U$, but from our framework,  $\hat{A}_U=(\zeta-1)\hat{p}_U$. Here we impose $\zeta=0$. 

\paragraph{Conclusion} Do we want to impose that $\zeta=0$?


\subsection{Back to square one?}

François, you've proposed to focus the paper on the methodology and the CCC scenarios. In this application, the main equation we use is $\hat{A}  = -\frac{\sum_i \Delta C^*_i}{GDP}$. Meaning that the aggregate productivity growth impact is the ratio of the extra net cost in 2050 and 2022 GDP. We find that the impact is -0.6\%, because the extra cost are 0.6\% of GDP.

It is actually pretty similar to the DICE-like approach. In \cite{dietzEndogenousGrowthConvexity2015}, GDP is 
\[
Y_t=(1-D_t^Y)(1-\Lambda_t)A_tK^{\alpha+\beta}_tL^{1-\alpha}_t
\]. $D_t$ is the damage to the economy, and $\Lambda_t$ is the emissions abatement costs as a share of GDP. $A_t$ is the exogenous element of TFP. One can express the endogenous TFP as $(1-\Lambda_t)A_t$. If we assume there is no growth of the exogenous element (or it is outside the scope of the analysis. Then $\hat{A}_t=\ln(1-\Lambda_t)\approx \Lambda_t= \frac{\sum_i \Delta C_{i,t,\tau}}{\text{GDP}}$. The impact of net zero on aggregate productivity is the sum of the net cost over GDP.

We are then back to what we were doing and did not satisfy you two years ago.

\paragraph{My take} We get back to this classic (and intuitive) result, only because we are implicitly making the assumption that $\zeta=0$. What we are adding is that if we don't make that simplifying assumption, then it just need a correction factor in front of it to take into account the dual nature of TFP (basically, the influence of $\hat{X}_i$).



\paragraph{We can show that we can dea}

Let's take the same framework as the section "Unprocessed vs unprocessed" and Eq. 4. And let's consider directly the aggregate productivity
\begin{equation*}
    \hat{A}_{t,\tau} =\lambda_M\hat{A}_{U,t,\tau}+\lambda_U\hat{A}_{U,t,\tau}
\end{equation*}
Using the dual definition of TFP growth
\begin{equation*}
\hat{A}_{t,\tau} =\lambda_M\left(-\hat{p}_M + \sum_j a_{jM} \hat{p}_{j} \right)+\lambda_U\left(-\hat{p}_U + \sum_j a_{jU} \hat{p}_{j} \right)
\end{equation*}
Multiplying each factor by $\frac{p_iX_i}{p_iX_i}$
\begin{equation*}
   \hat{A}_{t,\tau} = \frac{1}{p_MX_M}\lambda_M\left(-p_MX_M\hat{p}_M + \sum_j a_{jM}p_MX_M \hat{p}_{j} \right)+
     \lambda_U\frac{1}{p_UX_U}\left(-p_UX_U\hat{p}_U +\sum_j a_{jU} p_UX_U\hat{p}_{j}\right)
     \end{equation*}
Using $a_{ij}p_iX_i=X_{ji}$
\begin{equation*}
         \hat{A}_{t,\tau} = \frac{1}{p_MX_M}\lambda_M\left(-\Delta C_M + \sum_j X_{jM}\hat{p}_{j} \right)+
     \lambda_U\frac{1}{p_UX_U}\left(-\Delta C_U +\sum_j aX_{jU}\hat{p}_{j}  \right)
     \end{equation*}
Using $\Delta C_i^{U} = \sum_j X_{ji} \hat{p}_{j}$, and that no $U$'s input change prices.
\begin{equation*}
        \hat{A}_{t,\tau} = \frac{1}{p_MX_M}\lambda_M\left(-\Delta C_M + \Delta C_M^{U}\right)+
     \lambda_U\frac{1}{p_UX_U}\left(-\Delta C_U \right)
\end{equation*}
Reorganising, we get
\begin{align*}
    \hat{A}_{t,\tau} &= \frac{1}{p_MX_M}\lambda_M(-\Delta C_M)+  \left(\frac{1}{p_MX_M}\lambda_M+
    - \lambda_U\frac{1}{p_UX_U}\right)\Delta C_U \\
      &=- \frac{1}{GDP}\Delta C_M+  \left(\frac{1}{GDP}+
    - \frac{1}{GDP}\right)\Delta C_U \\
       &=- \frac{\Delta C_M}{GDP} \\
        &=- \frac{\Delta C_U+\Delta C_M^*}{GDP} \\
\end{align*}


In conclusion,
\begin{equation}
\label{eq_ahat_t_tau_ccc_no_zeta}
 \hat{A}_{t,\tau} = - \frac{\sum_i \Delta C^*_{i,t,\tau}}{\text{GDP}}.
\end{equation}

Two conclusions: 
\begin{enumerate}
    \item It shows that if we only consider the downstream sector using the full cost, then we capture all the productivity growth. If only one sector is included in the CCC, then processed and unprocessed cost are equal, and we are not missing anything by not considering the upstream sector. When the two sectors are in the CCC, then we should consider the processed cost.
    \item To arrive at this result, we do not need our framework at all. It stems from the dual definition of TFP. 
\end{enumerate}





\newpage



This result where we don't need $\zeta$ to express the industry-level (and then the aggregate in Eq. \ref{eq_ahat_processed_ccc}) stems from an implicit assumption that $\zeta=0$.

Under the assumption that $\hat{X}_i=0$, we transform the following fact (see appendix and eq 27):  $\hat{X}_i + \hat{p}_i = \frac{\hat{\Sigma}_p}{\zeta} = \frac{\hat{\Sigma}_X}{1-\zeta}=\hat{p}_i$. Using this we can express productivity growth three ways
\begin{equation*}
    \begin{cases}
    \text{Primal: } \hat{A}_i=\hat{X}_i - \hat{\Sigma}_{X_i}=-\hat{\Sigma}_{X_i}=-(1-\zeta) \hat{p}_i=(\zeta-1) \hat{p}_i\\
    \text{Dual: } \hat{A}_i=\hat{\Sigma}_{p_i} - \hat{p}_i=\zeta \hat{p}_i- \hat{p}_i=(\zeta-1) \hat{p}_i\\
       \text{Framework} \hat{A}_i=(\zeta-1)\hat{p}_i
\end{cases}
\end{equation*}

\ER{Actually, he's right, I cannot equal $-\frac{\Delta C^*_i}{p_i X_i}=(\zeta-1)\frac{\Delta C_i}{p_i X_i}$}
\er{Mais en fait ce qu'on dit c'est que pour chaque industry: 
\begin{equation*}
    \Delta C^*_i=(\zeta-1)\Delta C_i
\end{equation*}
}

Mais comme pour U, $  \Delta C^*_U=\Delta C_U$, alors $\zeta=0$. 

C'est vrai ça ? A checker dans les data du CCC. 


En revanche ça ça tiens:



\paragraph{Case 1}

A sector whose price change is only due to its own idiosyncratic productivity shock. Upstream sector. By assumption $\hat{X}_i=0$.

\begin{align*}
    \hat{A}_i &= -\left(\hat{p}_i - \sum_j a_{ji} \hat{p}_{j}\right)\\
    &=-\hat{p}_i
\end{align*}

But in the framework $ \hat{A}_i=(\zeta-1)\hat{p}_i$. 
Hence, to make it match, you implicitely imposes $\zeta=0$. 
It's still right to write: 
\begin{equation*}
    \Delta C^*_i=(\zeta-1)\Delta C_i
\end{equation*}

But also to write 

\begin{equation*}
   \hat{A}_i=(\zeta-1)\frac{\Delta C_i}{p_iX_i}
\end{equation*}

If you're saying that $\zeta$ is not calibrated on this specific sector, but is 0.126. Then $ \hat{A}_i=(\zeta-1)\hat{p}_i$ is wrong. 
We might know it in the case of the CCC (or assume we know it). But in general?

\paragraph{What is $\zeta$?}
You show that by writing the dual we get
\begin{equation*}
    \Delta C^*_i=(\zeta-1)\Delta C_i
\end{equation*}

So $\zeta$ takes care of the supply chain. That would be true if we were estimating 
\begin{equation*}
    \hat{A}_i =-\zeta_1\hat{p}_i+\zeta_2 \left(\sum_j a_{ji} \hat{p}_{j}\right)
\end{equation*}


\end{document}